\section{From particle configurations to the zeon quotient}
\label{hierarchy:collision-quotient}
% BEGIN SUBJECT INDEX: hierarchy:collision-quotient
\index{tensor power}
\index{symmetric power}
\index{collision!filtration}
\index{collision!free quotient}
\index{zeon algebra!collision-free interpretation}
\index{common-letter dynamics}
\index{independent-particle dynamics}
% END SUBJECT INDEX

Let $V$ have basis $e_1,\ldots,e_n$ over $\mathbb Q$. A deterministic
transformation $f$ sends $e_i$ to $e_{if}$ and acts on tensor powers by
\[
 (e_{i_1}\otimes\cdots\otimes e_{i_k})f^{\otimes k}
 =e_{i_1f}\otimes\cdots\otimes e_{i_kf}.
\]
This is simultaneous common-transformation motion of ordered particles.
Passing to $\operatorname{Sym}^k(V)$ forgets order but retains multiplicity:
repeated factors remain valid configurations.

\subsection{The invariant collision filtration}
Let $F_\ell$ be the span of degree-$k$ monomials with at most $\ell$
distinct occupied vertices. Deterministic transformations cannot increase
the number of distinct images, so
\[
 0\subseteq F_1\subseteq\cdots\subseteq F_k
 =\operatorname{Sym}^k(V)
\]
is an invariant filtration. The collision-free span itself need not be
invariant: distinct particles may collide. For $1\le k\le n$, the
quotient $F_k/F_{k-1}$ has basis the squarefree monomials $e_I$, indexed
by $k$-subsets, with action
\[
 e_I f^{\vee k}=\begin{cases}
 e_{If},&|If|=k,\\0,&|If|<k.
 \end{cases}
\]
This is the degree-$k$ zeon representation. Thus zeons isolate the
collision-free layer of the symmetric-power action rather than introduce
a different deterministic dynamics. Lost cardinality cannot be restored,
which proves $(fg)^{\vee k}=f^{\vee k}g^{\vee k}$.
An equivalent set-theoretic construction is the quotient of the module
on subsets of size at most $k$ by its submodule on subsets of size at most
$k-1$. This is not an exterior power: permutation signs are absent.

\subsection{Averages and limiting projections}
For deterministic generators $M_i$ chosen with probabilities $p_i$, put
\[
 A_k=\sum_i p_iM_i^{\vee k},\qquad
 \Omega_k=\lim_{m\to\infty}\frac1m\sum_{t=0}^{m-1}A_k^t.
\]
Each summand preserves the collision filtration before taking the quotient.
One must average these induced actions, not substitute the averaged
one-particle matrix into a symmetric-power formula. A general stochastic
linear map need not preserve the collision subspace of the symmetric
power. The independent-particle permanent construction is treated separately.

The basic observables remain visible in the full tensor construction.
Writing
\[
 \Phi(X)=\sum_i p_iM_iXM_i^{\mathsf T},\qquad
 \Psi(X)=\sum_i p_iM_i^{\mathsf T}XM_i,
\]
the limiting kernel measure gives
\[
 N=\lim_{m\to\infty}\frac1m\sum_{t=0}^{m-1}\Phi^t(I),\qquad
 M=\lim_{m\to\infty}\frac1m\sum_{t=0}^{m-1}\Psi^t(J).
\]
At zeon level two, $u_2=\Omega_2\mathbf1$ records survival of a distinct
pair, and its symmetric, zero-diagonal matrix realization is
\[
 \boxed{\operatorname{Mat}(u_2)=J-N.}
\]
The full projection $\Omega_2$ acts on pair space; $J-N$ is reconstructed
from its right field and is not that full projection. The $P_j,Q_i$
projections of Section~\ref{eproj:section} act on the original state space
and have yet another role. The following sections develop the higher
fields and their compatible inclusion maps within this framework.

